% =====================================================================
% Journée 4 --- S14 : TP CI/CD (pipeline GitHub Actions, service Avis-Clients)
% Compte rendu individuel --- ce qui a été fait, problèmes, solutions.
% Compilation : tectonic Devoir_J4_S14_CICD.tex
% =====================================================================
\documentclass[11pt,a4paper]{article}

\usepackage[utf8]{inputenc}
\usepackage[T1]{fontenc}
\usepackage[french]{babel}
\usepackage{mathptmx}            % Times New Roman
\usepackage[scaled=0.85]{helvet}
\usepackage{courier}
\usepackage[margin=2cm]{geometry}
\usepackage{booktabs}
\usepackage{tabularx}
\usepackage{ragged2e}
\usepackage{enumitem}
\usepackage{titlesec}
\usepackage{xcolor}
\usepackage{listings}
\usepackage{graphicx}
\usepackage{float}
\usepackage{caption}
\usepackage{fancyhdr}

\captionsetup{font=small,labelfont=bf,skip=4pt}
\graphicspath{{./}}

\definecolor{gristexte}{gray}{0.35}
\definecolor{codebg}{gray}{0.96}

\pagestyle{fancy}
\fancyhf{}
\fancyfoot[C]{\small\itshape\color{gristexte}Page \thepage}
\renewcommand{\headrulewidth}{0.4pt}
\renewcommand{\footrulewidth}{0pt}

\titleformat{\section}{\normalfont\Large\bfseries}{\thesection.}{0.5em}{}[\titlerule]
\titleformat{\subsection}{\normalfont\large\bfseries}{}{0em}{}
\titlespacing*{\section}{0pt}{14pt}{8pt}
\titlespacing*{\subsection}{0pt}{10pt}{4pt}

\setlist[itemize]{label=\textbullet,leftmargin=1.4em,itemsep=2pt,topsep=3pt}

\lstset{
  basicstyle=\ttfamily\small,
  backgroundcolor=\color{codebg},
  frame=none,
  xleftmargin=4pt,
  breaklines=true,
  showstringspaces=false,
  columns=fullflexible,
}

\newcolumntype{Y}{>{\RaggedRight\arraybackslash}X}
\newcolumntype{P}[1]{>{\RaggedRight\arraybackslash\hyphenpenalty=10000\exhyphenpenalty=10000}p{#1}}

\setlength{\parindent}{0pt}
\setlength{\parskip}{4pt}
\setlength{\emergencystretch}{2.5em}

% Capture : affiche l'image si présente, sinon un encadré "à insérer" (le doc compile quand même).
\newcommand{\shot}[2]{%
  \begin{figure}[H]\centering
  \IfFileExists{#1}{\includegraphics[width=\textwidth]{#1}}%
  {\fbox{\parbox[c][2.6cm][c]{0.95\textwidth}{\centering\itshape Capture à insérer : \texttt{#1}}}}%
  \caption{#2}\end{figure}}

\begin{document}

% ---------------------------------------------------------------------
\thispagestyle{empty}
\begin{center}
  {\LARGE\bfseries Journée 4}\\[4pt]
  {\large\bfseries Architecture Micro-services}\\[6pt]
  {\normalsize\itshape S14 \textperiodcentered{} TP CI/CD : pipeline automatisé du service Avis-Clients}\\[6pt]
\end{center}

\vspace{4pt}\hrule\vspace{8pt}

Nom / Prénom : DANILO Elouan \hfill Service : \textbf{Avis-Clients}

\vspace{4pt}

Dépôt : \texttt{https://github.com/Oasio/avis-clients-cicd}

Image : \texttt{ghcr.io/oasio/avis-clients-cicd:latest}

\vspace{6pt}

\section{Contexte}

Le TP consiste à doter mon service \textbf{Avis-Clients} (Symfony 7 + API Platform) d'un
\textbf{pipeline CI/CD complet}~: à chaque \texttt{push}, on lance les tests, on construit une
image Docker, et on la publie dans un \emph{Container Registry}, prête à déployer.

Le polycopié s'appuie sur \textbf{GitLab.com}. J'ai réalisé le TP sur \textbf{GitHub}, l'équivalent
présenté en cours (S13.1)~: \textbf{GitHub Actions} pour le pipeline (fichier
\texttt{.github/workflows/ci.yml}) et le \textbf{GitHub Container Registry} (\texttt{ghcr.io}) pour
l'image. La logique reste identique~: stages séquentiels, tag par hash de commit, déploiement
manuel. Les principales adaptations~: le runner GitHub fournit déjà Docker (pas besoin de
\emph{docker-in-docker}), et le \texttt{when: manual} de GitLab est reproduit par un
\textbf{environnement protégé}.

\section{Le pipeline mis en place}

Le workflow se déclenche \textbf{automatiquement à chaque \texttt{push} sur \texttt{main}}
(et manuellement via \texttt{workflow\_dispatch}). Il enchaîne quatre stages~:

\begin{itemize}
  \item \textbf{test} : PHPUnit sur une base SQLite jetable~;
  \item \textbf{build} : construction de l'image Docker à partir du \texttt{Dockerfile}~;
  \item \textbf{package} : push de l'image sur \texttt{ghcr.io} (tags \texttt{:<sha>} et \texttt{:latest})~;
  \item \textbf{deploy} \emph{(bonus)} : déploiement simulé, déclenché manuellement.
\end{itemize}

\subsection{Stage TEST}

Les tests d'API persistent réellement des avis~: il faut donc une base. Plutôt que de lancer un
serveur MariaDB dans la CI, je surcharge \texttt{DATABASE\_URL} vers une base \textbf{SQLite}
jetable, et je crée le schéma depuis les métadonnées de l'entité.

\begin{lstlisting}
test:
  runs-on: ubuntu-latest
  container: { image: php:8.3-cli }
  env:
    APP_ENV: test
    DATABASE_URL: "sqlite:///%kernel.project_dir%/var/data_test.db"
  steps:
    - uses: actions/checkout@v4
    - run: docker-php-ext-install intl pdo_mysql pdo_sqlite zip
    - run: composer install --no-interaction --prefer-dist
    - run: php bin/console doctrine:schema:create --no-interaction
    - run: php bin/phpunit          # 8 tests au vert
\end{lstlisting}

\subsection{Stages BUILD et PACKAGE}

Le stage \texttt{build} vérifie que l'image se construit. Le stage \texttt{package} se connecte au
registry avec le \texttt{GITHUB\_TOKEN} auto-généré (aucun identifiant en dur) puis pousse l'image
taguée par le hash court du commit \emph{et} par \texttt{latest}.

\begin{lstlisting}
package:
  needs: build
  permissions: { packages: write }
  steps:
    - run: echo "$GITHUB_TOKEN" | docker login ghcr.io -u $ACTOR --password-stdin
    - run: |
        IMAGE=$(echo "$IMAGE_NAME" | tr '[:upper:]' '[:lower:]')   # ghcr = minuscules
        docker build -t "$IMAGE:${GITHUB_SHA::7}" -t "$IMAGE:latest" .
        docker push "$IMAGE:${GITHUB_SHA::7}"
        docker push "$IMAGE:latest"
\end{lstlisting}

\subsection{Stage DEPLOY (bonus)}

GitHub Actions n'a pas de \texttt{when: manual} par job. Je reproduis le comportement avec un
\textbf{environnement \texttt{production}} protégé par un \emph{required reviewer}~: le job
\texttt{deploy} reste en attente tant qu'il n'est pas approuvé dans l'interface. Le déploiement
lui-même est simulé (un \texttt{echo}), aucun serveur n'étant requis pour le TP.

\section{Problèmes rencontrés et solutions}

\begin{tabularx}{\textwidth}{@{}P{6.2cm} Y@{}}
\toprule
\textbf{Problème} & \textbf{Solution mise en place} \\
\midrule
Les tests d'API (\texttt{WebTestCase}) écrivent en base~; aucun serveur MariaDB n'est disponible
dans le runner CI. &
Surcharge de \texttt{DATABASE\_URL} vers une base \textbf{SQLite} jetable dans le job. Les variables
d'environnement réelles étant prioritaires sur les fichiers \texttt{.env}, la config MariaDB du
projet reste intacte. \\[4pt]
La migration livrée (\texttt{Version20260612000000}) est du SQL \textbf{MariaDB pur}
(\texttt{AUTO\_INCREMENT}, \texttt{ENGINE=InnoDB}), non rejouable sur SQLite. &
En CI, création du schéma via \texttt{doctrine:schema:create} (généré depuis les métadonnées de
l'entité, indépendant du SGBD) au lieu de rejouer la migration. \\[4pt]
Le registry \texttt{ghcr.io} refuse les noms d'image en majuscules, or
\texttt{github.repository} contient \texttt{Oasio}. &
Passage en minuscules dans le script avec \texttt{tr '[:upper:]' '[:lower:]'} avant chaque
\texttt{docker build/push}. \\[4pt]
Authentification au registry sans écrire de secret. &
Login via le \texttt{GITHUB\_TOKEN} auto-généré et permission \texttt{packages: write} sur le job,
soit l'équivalent des variables \texttt{\$CI\_REGISTRY\_*} de GitLab. \\[4pt]
Reproduire le \texttt{when: manual} de GitLab, absent de GitHub Actions. &
Job \texttt{deploy} rattaché à un \textbf{environnement \texttt{production}} protégé par un
\emph{required reviewer}~: il ne démarre qu'après approbation manuelle. \\
\bottomrule
\end{tabularx}

\section{Résultat}

Le dernier pipeline passe \textbf{intégralement au vert} (statut \emph{Success}, 5\,min\,15\,s),
et l'image est bien publiée dans le registry, taguée \texttt{latest} et par le hash du commit.
Le déclenchement est automatique à chaque \texttt{push} sur \texttt{main}, et tous les stages
attendus (test, build, package) sont exécutés, plus le déploiement manuel en bonus.

\section{Conclusion}

Mon service Avis-Clients dispose désormais d'un pipeline CI/CD versionné qui, à chaque \texttt{push},
teste le code, construit l'image Docker et la publie sur \texttt{ghcr.io}, prête à déployer. Les
deux points qui m'ont demandé le plus d'ajustement, faire tourner les tests sans serveur de base
et publier proprement sur le registry, ont été réglés respectivement par une base SQLite jetable
(schéma créé depuis l'entité) et par un login au registry via le jeton auto-généré. Le déploiement
manuel via environnement protégé apporte le point bonus tout en évitant tout déploiement accidentel.

% ---------------------------------------------------------------------
\newpage
\appendix
\section{Annexes : captures d'écran}

\shot{s14-pipeline.png}{Annexe A. Dernier pipeline GitHub Actions en \textbf{succès} (\texttt{Success}, 5\,min\,15\,s). Les quatre stages \texttt{test} \textrightarrow{} \texttt{build} \textrightarrow{} \texttt{package} \textrightarrow{} \texttt{deploy} sont verts~; déclenché via \texttt{push} sur \texttt{main}.}
\shot{s14-registry.png}{Annexe B. GitHub Container Registry (\texttt{ghcr.io}) : l'image \texttt{avis-clients-cicd} est publiée et taguée \texttt{latest} (et par le hash court du commit). Commande de récupération~: \texttt{docker pull ghcr.io/oasio/avis-clients-cicd:latest}.}

\end{document}
